\documentclass[11pt,letterpaper]{article}

\usepackage[margin=1in]{geometry}
\usepackage[T1]{fontenc}
\usepackage{amsmath}
\usepackage{amssymb}
\usepackage{array}
\usepackage{booktabs}
\usepackage{graphicx}
\usepackage{xcolor}
\usepackage[numbers,sort&compress]{natbib}
\usepackage[colorlinks=true,linkcolor=blue!50!black,citecolor=blue!50!black,urlcolor=blue!50!black]{hyperref}

% Absorbs the occasional long compound word without visibly loosening paragraphs.
\setlength{\emergencystretch}{2em}

% Visible marker for facts and numbers that the experiments must still supply.
\newcommand{\ph}[1]{\textcolor{red}{[TODO: #1]}}
\newcommand{\efface}{\textsc{Efface}}

\title{\efface: Finding the Most Effective Meaning-Preserving Attack Chain\\against SynthID-Text-Class Watermarks}

\author{
  Guillaume Meyer \\
  \ph{affiliation} \\
  \texttt{guillaume.meyer@outlook.com}
}

\date{Draft --- \today}

\begin{document}

\maketitle

\begin{abstract}
Text watermarking schemes of the SynthID-Text class embed a statistical signal in the sampling process of a language model, so the mark lives in the choice of tokens rather than in attached metadata. A growing literature measures how well such marks survive individual transformations: paraphrasing, back-translation, token substitution, and regeneration. That literature evaluates one transformation at a time, on its own operating point, and rarely reports how much meaning the transformation costs. This paper asks a joint question. Given a fixed semantic-preservation floor, which single transformation, or which ordered chain of at most three transformations, removes the watermark most reliably? We formalise the problem as a constrained search over an attack catalogue, measure removal against a same-config detector under the generating key and at a corpus-level true-positive rate, and select chains with a document-level composite of BERTScore and COMET, with the other metrics reported beside it. The design evaluates the selection on a held-out split, clusters every interval by prompt, and audits the chains the search skips. We release the catalogue, the search procedure, and the evaluation harness so that the frontier between removal and preservation can be reproduced and extended.
\end{abstract}

\section{Introduction}
\label{sec:intro}

Provenance marks for generated text are now deployed at scale. SynthID-Text, the scheme Google ships with Gemini, biases the sampling distribution at generation time and detects the bias later from the token sequence alone \citep{synthidtext}. The same design space contains green-list schemes that boost a keyed subset of the vocabulary \citep{kgw}, exponential and Gumbel variants that key the sampler noise \citep{kuditipudi}, and distortion-free constructions with formal guarantees \citep{christ}. What these schemes share is that the signal is spread thinly across every token: a mark is a statistical property of the whole passage, and each sentence carries only a small part of the evidence.

That property makes removal a question of degree. Editing a heading, reordering paragraphs, or fixing a comma leaves the token-level signal almost untouched, as the reliability study of \citet{kgwreliability} shows. Removing a mark requires replacing a substantial share of the original word choices. Every replacement is also a change in meaning, tone, and specificity, and the change compounds when several transformations are applied in sequence. Practitioners therefore face a trade-off rather than a binary outcome, and the trade-off has not been mapped systematically.

Existing evaluations each cover a slice of the space. The reliability study of \citet{kgwreliability} measures machine and human paraphrasing and copy-paste mixing. DIPPER establishes that a strong discourse-level paraphraser can strip detection while keeping meaning close \citep{dipper}. The cross-lingual study of \citet{cwra} shows that a pivot-language round trip removes marks that survive same-language paraphrase. MarkMyWords compares a suite of perturbations across schemes and reports quality alongside detectability \citep{markmywords}. None of these works searches over chains, and none reports which combination is optimal under a stated semantic budget.

This paper closes that gap for the SynthID-Text class. We treat removal as a constrained optimisation: maximise the fraction of documents whose mark is no longer detectable, subject to a semantic-preservation score at or above a fixed floor, over single tactics and over ordered chains of up to three tactics. The constraint is what makes the question well posed. Without it, the strongest tactic is simply the most destructive one, which the impossibility results of \citet{sand} predict must exist.

\paragraph{Contributions.}
\begin{itemize}
  \item A catalogue of thirteen rows: twelve removal tactics in four families and one control, each defined operationally, linked to the paper that introduced or evaluated it, and implemented behind one interface (\S\ref{sec:catalogue}).
  \item A chain search that evaluates ordered compositions of up to three tactics under a semantic floor, with pruning that keeps the evaluation budget bounded (\S\ref{sec:search}).
  \item An empirical frontier between removal and preservation across \ph{number} documents, \ph{number} generators, and \ph{number} domains, reported at two operating points (\S\ref{sec:results}).
\end{itemize}

\paragraph{Findings.}
\ph{One paragraph once the runs finish: best single tactic, best pair, best triple, and the marginal value of the third step.}

\section{Related Work}
\label{sec:related}

\paragraph{Sampling-time text watermarks.}
Green-list watermarking partitions the vocabulary at each step with a keyed hash of recent context and biases sampling toward the green set; detection counts green tokens against a null \citep{kgw}. SynthID-Text generalises the same idea with a tournament sampling procedure designed for production deployment and for detection from text alone \citep{synthidtext}. Exponential and Gumbel schemes key the sampler's noise and admit exact tail tests \citep{kuditipudi}. Distortion-free schemes sample from the original distribution while retaining a detectable coupling \citep{christ}. Our study targets the SynthID-Text class because it is deployed, because an open reference implementation exists, and because the detection test is well specified; we expect the qualitative ranking of tactics to transfer to other sampling-bias schemes, and we treat that transfer as a question for future work.

\paragraph{Removal attacks.}
Paraphrasing is the canonical attack. \citet{kgwreliability} show that machine paraphrase reduces detection at short lengths and that the mark recovers as the passage grows, and that human rewriting is stronger than machine rewriting. DIPPER, an 11B discourse-level paraphraser, evades several detectors including green-list watermarking while preserving meaning closely \citep{dipper}. Token-substitution attacks replace a fraction of words with synonyms and were used as the robustness probe in \citet{threebricks}. A more targeted line masks high-self-information tokens and infills them, which is the masked-LM infilling row of our catalogue \citep{sira}. Regeneration with a different model replaces the token stream wholesale. The impossibility result of \citet{sand} gives a quality-oracle random walk that removes any watermark that occupies a small share of the high-quality output space, and \citet{sadasivan} argue more broadly that detection is fragile under paraphrase. \citet{jovanovic} study the dual problem of stealing a scheme and forging marks, show that stealing also boosts scrubbing, and bound what a detector can be trusted to report. An independent evaluation applies the attack to SynthID-Text and reports that an off-the-shelf paraphraser clears the mark in over 90\% of roughly 1,000-token samples at a \(10^{-3}\) false-positive rate \citep{probingsynthid}.

\paragraph{Back-translation and cross-lingual removal.}
Round-trip translation through a pivot language perturbs token choice while preserving propositional content, which makes it an attractive removal primitive. \citet{cwra} formalise the cross-lingual removal attack, show that watermark strength is not consistent across languages, and report a large drop in detection after a pivot round trip. MarkMyWords includes translation among its tamper-resistance perturbations and reports that long generations often remain detectable \citep{markmywords}. The two results differ in scheme, length, and pivot, which is one reason a joint evaluation under a common operating point is useful.

\paragraph{Benchmarks and tooling.}
MarkMyWords unifies several schemes behind one interface and reports quality, detection length, and tamper resistance \citep{markmywords}. MarkLLM provides reference implementations of the KGW and SynthID scheme families together with generation, detection, and visualisation utilities \citep{markllm}. We build on MarkLLM for the detector and keep our harness independent of any single scheme so that the catalogue and the search can be reused.

\paragraph{Evaluation of preservation.}
Surface overlap such as BLEU is a poor proxy for meaning under paraphrase. We therefore combine embedding-based similarity \citep{bertscore}, learned reference-based quality \citep{bleurt}, learned translation quality \citep{comet}, and distributional closeness \citep{mauve}. Section~\ref{sec:metrics} states how the signals are aggregated.

\section{Problem Formulation}
\label{sec:problem}

\paragraph{Setup.}
A generator $G$ produces a document $x = (x_1, \dots, x_n)$ under a watermark key $k$ drawn from a keyspace $\mathcal{K}$. A detector $D$ takes a document and a key and returns a score whose null distribution is known under the unwatermarked hypothesis; we use the $p$-value $p_k(x)$ reported by the same-config test, so that small values indicate the presence of the mark. The detector is valid only against the scheme configuration and key family used at generation, a limitation we return to in \S\ref{sec:limitations}.

\paragraph{Attack catalogue.}
Let $\mathcal{A} = \{t_1, \dots, t_m\}$ be a catalogue of transformations. Each $t$ maps a document to a document, $t: \mathcal{X} \to \mathcal{X}$, and carries an intensity parameter $\theta_t$ that controls how much it changes. Table~\ref{tab:catalogue} lists the catalogue. A transformation may be stochastic, in which case $t(x)$ denotes a draw from its output distribution.

\paragraph{Chains.}
An attack chain is an ordered tuple $T = (t^{(1)}, \dots, t^{(K)})$ with $K \le 3$, applied left to right:
\begin{equation}
  x^{(0)} = x, \quad x^{(j)} = t^{(j)}\!\left(x^{(j-1)};\, \theta^{(j)}\right), \quad x^{(K)} = T(x).
\end{equation}
Order matters. A paraphrase followed by a round trip yields a different document than a round trip followed by a paraphrase, because each step conditions on the output of the previous one. We therefore treat chains as ordered and search over ordered tuples.

\paragraph{Removal.}
For a document $x$ and a key set $K = \{k_1, \dots, k_L\}$ of independent keys, define the document as removed when the mark fails to be detected under every key at level $\alpha$:
\begin{equation}
  \mathrm{removed}_\alpha(x) = \mathbf{1}\!\left[\min_{k \in K} p_k(x) > \alpha \right].
\end{equation}
Two summaries follow. The \emph{removal rate} is the mean of $\mathrm{removed}_\alpha$ over the evaluation set, reported with a Wilson interval. The \emph{corpus operating point} is the true-positive rate at a $1\%$ false-positive rate. Positives are watermarked documents, attacked or not; negatives are unwatermarked documents. The threshold is the highest detector score at which the negative false-positive rate does not exceed $1\%$. We report both, because a per-document threshold and a corpus-level rate answer different questions: the first asks whether this document still carries a mark, the second asks what a detector tuned for a fixed error budget would conclude.

\paragraph{Preservation.}
Let $S(x, T(x))$ be a semantic-preservation score in $[0,1]$, aggregated from the metrics in \S\ref{sec:metrics}. The floor $\tau$ is fixed before the search and reported with every result. A chain that falls below $\tau$ is inadmissible, whatever its removal rate.

\paragraph{Objective.}
The problem we study is
\begin{equation}
  T^\star = \operatorname*{arg\,max}_{T \in \mathcal{T}_{\le 3}} \; \mathbb{E}\left[\mathrm{removed}_\alpha(T(x))\right]
  \quad \text{subject to} \quad
  \mathbb{E}_{x}\left[S(x, T(x))\right] \ge \tau ,
  \label{eq:objective}
\end{equation}
where $\mathcal{T}_{\le 3}$ is the set of ordered chains of length at most three drawn from $\mathcal{A}$. The constraint is enforced on the mean score; for every chain we recommend we also report the fraction of documents that fall below $\tau$. Equation~\ref{eq:objective} makes the trade-off explicit and prevents the degenerate answer, which is to apply the most destructive transformation available.

\paragraph{Why a floor is necessary.}
\citet{sand} show that for any watermark occupying a small share of the space of high-quality responses, a quality-oracle random walk reaches an unmarked output. The construction needs a quality oracle and enough mixing steps, and it degrades quality to the extent that the walk is imperfect. A study that reports removal without a preservation constraint therefore measures the quality of the attacker's oracle as much as the robustness of the scheme. Fixing $\tau$ and reporting the frontier around it keeps the comparison meaningful.

\section{Attack Catalogue}
\label{sec:catalogue}

Table~\ref{tab:catalogue} lists the tactics we study. Each row is a transformation with one interface, one intensity parameter, and one implementation in the repository. The systems we execute are prompted rewrites, a local masked-LM filler, and deterministic text rules; several rows are operationalisations of a published attack rather than the published system. The row-by-row difference is in the released catalogue document: discourse and document rewrite prompt a model rather than running DIPPER, masked infilling masks a seeded sample rather than the highest-self-information tokens, synonym substitution takes candidates from a prompted model rather than Three Bricks' substitution model, and the pivot rows use one model for every leg. The search is non-adaptive: no tactic is trained against the detector. The families differ in what they perturb: paraphrase and regeneration replace the token stream through a model, back-translation replaces it through a translation system, and token edits touch individual positions. The control row is included so that a reader can see the difference between removing an edit-based carrier and attacking a statistical mark.

\begin{table*}[t]
\centering
\footnotesize
\caption{Attack catalogue. Cost is the expected loss of semantic fidelity at the intensity used in the main experiments: low, medium, or high. Source names the closest published relative, not necessarily the system executed; \S\ref{sec:catalogue} states the operational difference per row. Each row maps to a module under \texttt{efface/attacks/}.}
\label{tab:catalogue}
\begin{tabular}{@{}>{\raggedright\arraybackslash}p{3.4cm} >{\raggedright\arraybackslash}p{2.6cm} >{\raggedright\arraybackslash}p{3.5cm} l >{\raggedright\arraybackslash}p{3.3cm}@{}}
\toprule
Tactic & Family & What it changes & Cost & Closest relative \\
\midrule
Instruction paraphrase      & Paraphrase      & Sentence wording, kept close to source & Low    & \citet{kgwreliability} \\
Discourse paraphrase        & Paraphrase      & Paragraph structure and wording & Medium & \citet{dipper} \\
Masked-LM infilling         & Paraphrase      & Individual spans, local context only & Low & \citet{sira} \\
Style normalisation         & Paraphrase      & Contractions, casing, punctuation, filler & Low & \citet{markmywords} \\
Document rewrite            & Paraphrase      & Whole document, fresh context per chunk & High & \citet{dipper} \\
Single-pivot back-translation & Back-translation & Token stream via one pivot language & Medium & \citet{markmywords} \\
Multi-hop back-translation  & Back-translation & Token stream via two or more pivots & High & \citet{cwra} \\
Cross-lingual generation    & Back-translation & Response generated in the pivot, then translated & High & \citet{cwra} \\
Synonym substitution        & Token edits     & A fraction of content words & Low & \citet{threebricks} \\
Random word edits           & Token edits     & Swap, insert, or delete at random positions & Low & \citet{markmywords} \\
Typo and noise injection    & Token edits     & Character-level noise and misspellings & Low & \citet{markmywords} \\
Non-origin regeneration     & Regeneration    & The entire token stream, from a different model & High & \citet{cwra} \\
\midrule
Unicode scrub (control)     & Control         & Invisible and zero-width characters only & None & \citet{markmywords} \\
\bottomrule
\end{tabular}
\end{table*}

\paragraph{Paraphrase.}
Paraphrase replaces the writer's word choices while keeping the propositional content. The family spans a range of intensity. A style pass changes surface form and leaves most content words in place. Instruction-prompted paraphrase rewrites sentence by sentence and typically keeps most entities and numbers. DIPPER-style discourse paraphrase conditions on the surrounding passage and may reorder or merge sentences \citep{dipper}. Masked-LM infilling replaces a fraction of spans with predictions from a masked language model, which keeps the rest of the token stream intact \citep{sira}. Document rewriting discards local context and can drift furthest from the source, which is why it sits at the expensive end.

\paragraph{Back-translation.}
A pivot round trip translates the document into a pivot language and back. The output is fluent in the source language and usually preserves propositions, while the token sequence is the translator's rather than the generator's. \citet{cwra} show that watermark strength is not preserved across languages and that the round trip removes a large share of the signal. Multi-hop chains repeat the round trip through several pivots. Cross-lingual generation differs in what happens in the middle of the chain: the source is translated into the pivot, the attacker model rewrites the translated document in the pivot language, and the rewrite is translated back. That variant moves the watermark closer to its source of failure, because the middle step replaces the token stream in a language the watermark was never applied to.

\paragraph{Token edits.}
Token edits touch a bounded fraction of positions. Synonym substitution preserves meaning to the extent that the substitution model is reliable \citep{threebricks}. Random swaps, insertions, and deletions are weaker per edit but cheap to apply at scale \citep{markmywords}. Typo and noise injection perturbs characters, which changes tokenisation and therefore the context used to key the watermark. These tactics are the cheapest per document. A later full-stream replacement overwrites them, so the search treats them as last steps rather than as prefixes a rewrite then discards (\S\ref{sec:search}).

\paragraph{Regeneration.}
Regeneration asks a different model to produce a new document from the prompt, with the original available or withheld. It replaces every token and therefore removes the mark most completely, at the highest cost to fidelity. We include it as the upper bound on removal and the lower bound on preservation.

\paragraph{Control.}
The Unicode scrub removes zero-width characters, bidirectional controls, and exotic spaces. These carriers belong to edit-based watermarking, where the mark is inserted into the text as characters. The control row quantifies how much of an observed removal is attributable to that channel, and it should show no effect on a sampling-time mark.

\section{Searching over Chains}
\label{sec:search}

\paragraph{Space size.}
With $m = 12$ tactics and chains of length up to three, the ordered search space has $m + m^2 + m^3 = 1{,}884$ chains before intensities are considered. Each chain requires a generation-free pass over \ph{number} documents, a detector evaluation under $L$ keys, and the semantic metrics. Two intensities per step would multiply the space by eight. An exhaustive sweep is possible in principle and too expensive in practice, so we search.

\paragraph{Procedure.}
We use a beam search over depth. At depth one, every tactic is evaluated at its default intensity and scored by the objective in Equation~\ref{eq:objective}. At depth two, each surviving chain is extended with every tactic, including repeats, and re-evaluated. At depth three the same extension applies, truncated to the best $B$ chains. The beam width $B$ is fixed before the search and reported. Every chain is scored on the same document set, so comparisons within a depth are paired. Across depths the comparison is descriptive, since the admissible set changes.

\paragraph{Pruning.}
Two rules keep the budget bounded. A chain whose mean preservation score falls below $\tau - \epsilon$ is dropped, where $\epsilon$ accounts for metric noise. A chain whose removal rate does not exceed that of its best prefix by more than the smallest detectable difference is dropped, because the extra step costs meaning without buying removal. Both rules are logged, so a reader can see which chains were evaluated and which were inferred to be dominated.

\paragraph{Intensity.}
Each tactic has one intensity parameter, normalised to $[0,1]$ where the implementation permits. The main sweep uses \ph{number} settings per tactic. For back-translation the intensity is the number of pivots; for token edits it is the fraction of positions; for paraphrase it is a model-side control such as lexical diversity. We report the chosen intensity with every recommended chain, because a chain that only works at maximum intensity is a different recommendation from one that works at a moderate setting.

\paragraph{Selection and multiple comparisons.}
Searching 1{,}884 chains and reporting the best invites selection bias. We therefore split documents into a search split and a held-out split before any evaluation, choose the chain on the search split, and report its performance on the held-out split. The headline numbers in \S\ref{sec:results} are held-out numbers. Where the budget does not allow a full held-out pass for every chain, we report the top \ph{number} chains on both splits and say so explicitly.

\paragraph{Stopping.}
The search stops at depth three by construction, and earlier when no chain survives pruning. We also report the marginal effect of the third step, which is the quantity that decides whether the three-tactic budget is worth spending at all. If the best pair captures most of the achievable removal, the honest recommendation is a pair.

\section{Experimental Setup}
\label{sec:setup}

\paragraph{Generators and schemes.}
Table~\ref{tab:generators} lists the generators. Both rows are open-weight models, so the sampling distribution is available and the same SynthID reimplementation can be applied to each. The G1 row uses \texttt{facebook/opt-1.3b}, the default model of the MarkLLM SynthID configuration. The G2 row uses \texttt{Qwen/Qwen2.5-1.5B-Instruct}, from a different family. Neither row is hosted: a hosted API exposes no logits, and Google retired text-watermark detection on its API in August 2026. The detector is MarkLLM's SynthID reimplementation, a research implementation with its own keying, and every removal number is relative to it.

\begin{table}[t]
\centering
\small
\caption{Generators and the scheme each row carries. Logits indicates whether the sampling distribution is available to the watermarking implementation.}
\label{tab:generators}
\begin{tabular}{@{}llll@{}}
\toprule
Row & Model & Scheme & Logits \\
\midrule
G1 & \texttt{facebook/opt-1.3b} & MarkLLM SynthID reimplementation & Yes \\
G2 & \texttt{Qwen/Qwen2.5-1.5B-Instruct} & MarkLLM SynthID reimplementation & Yes \\
\bottomrule
\end{tabular}
\end{table}

\paragraph{Attackers.}
Every tactic that calls a model uses a model from a family different from the generator, so that the attacker cannot reintroduce the generator's own mark. Table~\ref{tab:attackers} lists the models and their roles. The rule is applied per generator row: a chain that attacks a G1 document never uses a G1-family model, and the same holds for G2.

One attacker is used per chain, so a chain's result is attributable to the tactic sequence rather than to a model switch inside the sequence. The search and the intensity sweep use DeepSeek Flash. After the search, the top five chains, the best single tactic, and the Unicode control are re-run end to end with Claude Sonnet on the held-out split. That ablation tests whether the ranking transfers across attacker capability, and its cost stays a small fraction of the search. Exact model snapshots and prices are pinned in \texttt{configs/attackers.yaml} and recorded in every run manifest.

\begin{table}[t]
\centering
\small
\caption{Models used by the attack tactics.}
\label{tab:attackers}
\begin{tabular}{@{}>{\raggedright\arraybackslash}p{3.4cm} >{\raggedright\arraybackslash}p{3.2cm} >{\raggedright\arraybackslash}p{7.2cm}@{}}
\toprule
Role & Model & Used by \\
\midrule
Paraphrase and style     & DeepSeek Flash & Instruction paraphrase, discourse paraphrase, style, document rewrite \\
Translation              & DeepSeek Flash & All back-translation tactics \\
Masked infilling         & \texttt{roberta-large} & Masked-LM infilling \\
Synonym candidates       & DeepSeek Flash & Synonym substitution \\
Regeneration             & DeepSeek Flash & Non-origin regeneration \\
Second-attacker ablation & Claude Sonnet & Top five chains, best single, control \\
\bottomrule
\end{tabular}
\end{table}

\paragraph{Data.}
Table~\ref{tab:data} lists the prompt sets. Three domains cover distinct register and length profiles, and the MarkMyWords tasks are included so that at least part of the evaluation can be compared with published tamper-resistance numbers. Each domain contributes 100 prompts. With four domains and three keys, each generator row contributes 1{,}200 documents, and the two rows contribute 2{,}400 before exclusions.

\begin{table}[t]
\centering
\small
\caption{Prompt sets.}
\label{tab:data}
\begin{tabular}{@{}llll@{}}
\toprule
Domain & Source & Prompts & Target length \\
\midrule
News writing        & \ph{source} & 100 & \ph{tokens} \\
Long-form question answering & \ph{source} & 100 & \ph{tokens} \\
Creative writing    & \ph{source} & 100 & \ph{tokens} \\
MarkMyWords tasks   & \citet{markmywords} & 100 & \ph{tokens} \\
\bottomrule
\end{tabular}
\end{table}

\paragraph{Generation settings.}
Decoding uses temperature $0.7$, top-$p$ $0.9$, and $400$ new tokens for every generator, with the same prompt set across rows. Each prompt is generated once per key. Documents below 150 tokens are discarded, because the detector's power depends on length and short documents would be removed trivially. We report the number discarded.

\paragraph{Keys and repetitions.}
Each document is watermarked under $L = 3$ independent keys drawn from the scheme's keyspace and never reused across runs. Detection is evaluated under the key used at generation. Where the tactic is stochastic, each document passes through the tactic three times with different seeds, and the reported score is the mean. Seeds are fixed and recorded in the configuration files.

\paragraph{Precision and power.}
The $60/40$ split leaves 720 documents per generator row for the search and 480 for the held-out report. At a removal rate near $0.5$, a $95\%$ Wilson interval spans about $\pm 4.5$ percentage points on 480 documents and $\pm 3.6$ on 720; a per-domain cell on the held-out split has 120 documents and spans about $\pm 8.9$ points. Averaging three seeds over a binary outcome puts each document's score on a $1/3$ grid, which sets the per-document resolution; the interval on the mean is set by the document count. The preregistered minimum detectable difference is $0.05$ absolute removal rate. A paired comparison of two chains on the same 480 documents reaches $80\%$ power at that difference when the within-document correlation between the two chains' outcomes is about $0.7$ or higher, and less when it is lower; the pilot run estimates the correlation, and we report the shortfall if the design misses the target. A fallback design with two domains at 60 prompts would leave 144 held-out documents per row; there the same difference is below $50\%$ power even at a within-pair correlation of $0.8$.

\paragraph{Oracle.}
Detection uses the SynthID scheme implementation in MarkLLM \citep{markllm}, configured with the same parameters used at generation. The per-document $p$-value comes from a normal approximation to the unwatermarked null, with variance $1/(12\,n_{\mathrm{eff}})$ over the $n_{\mathrm{eff}}$ unmasked positions; we calibrate it on 100 unwatermarked documents per generator row. This is a same-config oracle: it detects the scheme it was given, under the key family it was given. It does not certify what a vendor detector would report, and we make no such claim. At the corpus operating point, 100 negatives admit one false positive, so that point is coarser than the per-document rate. Section~\ref{sec:limitations} states the consequences.

\paragraph{Compute.}
All runs complete on \ph{hardware description} in \ph{wall-clock time}. The attacker-model calls account for \ph{share} of the cost, and the detector passes for \ph{share}. Appendix~\ref{app:hyper} lists the per-tactic settings, and the configuration files in \texttt{configs/} are the authoritative record.

\section{Results}
\label{sec:results}

\subsection{Metrics and Decision Rule}
\label{sec:metrics}

\paragraph{Removal.}
We report the removal rate at $\alpha = 0.05$ under the generating key, with a $95\%$ prompt-clustered bootstrap interval and a document-level Wilson interval beside it, and the corpus-level true-positive rate at a $1\%$ false-positive rate from the parametric null. The first is the per-document question, the second is the detector-tuning question, and they can disagree: a chain can push most documents below the per-document threshold while leaving a detectable residual in the pooled distribution. Reporting both makes the disagreement visible.

\paragraph{Preservation.}
Two signals form the selecting composite. BERTScore-F1 measures embedding-level similarity to the source \citep{bertscore}. COMET measures adequacy and fluency, with the source document on the source side, the attacked document as the hypothesis, and the watermarked source as the reference \citep{comet}. Both are normalised to $[0,1]$, and $S$ is their unweighted mean. BLEURT provides a second learned judgement \citep{bleurt} and MAUVE measures distributional closeness between the attacked and unattacked sets \citep{mauve}; both are reported on the same documents and do not enter $S$ or $\tau$, so adding BLEURT at the larger tier cannot move which chain is selected. We also report the length ratio and the achieved edit rate, because a tactic that shortens the document can inflate similarity scores while dropping content. Per-metric values are always reported so a reader can reweight them, and the fraction of documents below $\tau$ is reported for every recommended chain.

\paragraph{Threshold.}
The semantic floor is $\tau = 0.80$, fixed before the search and revised at most once after the pilot run, with any revision recorded as a dated deviation in the preregistration. Results at $\tau - 0.05$ and $\tau + 0.05$ are reported in Appendix~\ref{app:perdomain} so that the sensitivity of the recommendation to the floor is visible. The attack prompts are in Appendix~\ref{app:prompts}.

\subsection{Single Tactics}

Table~\ref{tab:singles} reports every tactic applied once. The control row is the reference point: its removal rate should match the unattacked baseline.

\begin{table}[t]
\centering
\small
\caption{Single tactics at their default intensity. Higher removal is better; higher TPR at $1\%$ FPR means the detector still finds the mark. Preservation columns are higher-is-better except length ratio, which is closest to $1$ at $1.0$.}
\label{tab:singles}
\setlength{\tabcolsep}{4pt}
\begin{tabular}{@{}>{\raggedright\arraybackslash}p{4.0cm}ccccc@{}}
\toprule
Tactic & Removal $\uparrow$ & Residual TPR@1\%FPR $\downarrow$ & BERTScore $\uparrow$ & COMET $\uparrow$ & Length ratio \\
\midrule
Instruction paraphrase   & \ph{run} & \ph{run} & \ph{run} & \ph{run} & \ph{run} \\
Discourse paraphrase     & \ph{run} & \ph{run} & \ph{run} & \ph{run} & \ph{run} \\
Masked-LM infilling      & \ph{run} & \ph{run} & \ph{run} & \ph{run} & \ph{run} \\
Style normalisation      & \ph{run} & \ph{run} & \ph{run} & \ph{run} & \ph{run} \\
Document rewrite         & \ph{run} & \ph{run} & \ph{run} & \ph{run} & \ph{run} \\
Single-pivot back-translation & \ph{run} & \ph{run} & \ph{run} & \ph{run} & \ph{run} \\
Multi-hop back-translation & \ph{run} & \ph{run} & \ph{run} & \ph{run} & \ph{run} \\
Cross-lingual generation & \ph{run} & \ph{run} & \ph{run} & \ph{run} & \ph{run} \\
Synonym substitution     & \ph{run} & \ph{run} & \ph{run} & \ph{run} & \ph{run} \\
Random word edits        & \ph{run} & \ph{run} & \ph{run} & \ph{run} & \ph{run} \\
Typo and noise injection & \ph{run} & \ph{run} & \ph{run} & \ph{run} & \ph{run} \\
Non-origin regeneration  & \ph{run} & \ph{run} & \ph{run} & \ph{run} & \ph{run} \\
\midrule
Unicode scrub (control)  & \ph{run} & \ph{run} & \ph{run} & \ph{run} & \ph{run} \\
\bottomrule
\end{tabular}
\end{table}

\subsection{Pairs and Triples}

Tables~\ref{tab:pairs} and~\ref{tab:triples} report the best chains at each depth on the held-out split, ordered by removal rate among chains that satisfy the semantic floor. The full ranked list is released with the results archive.

\begin{table}[t]
\centering
\small
\caption{Best ordered pairs. Ranks are by removal rate among admissible chains.}
\label{tab:pairs}
\begin{tabular}{@{}clcccc@{}}
\toprule
Rank & Chain & Removal $\uparrow$ & Residual TPR@1\%FPR $\downarrow$ & COMET $\uparrow$ & Length ratio \\
\midrule
1 & \ph{chain} & \ph{run} & \ph{run} & \ph{run} & \ph{run} \\
2 & \ph{chain} & \ph{run} & \ph{run} & \ph{run} & \ph{run} \\
3 & \ph{chain} & \ph{run} & \ph{run} & \ph{run} & \ph{run} \\
4 & \ph{chain} & \ph{run} & \ph{run} & \ph{run} & \ph{run} \\
5 & \ph{chain} & \ph{run} & \ph{run} & \ph{run} & \ph{run} \\
\bottomrule
\end{tabular}
\end{table}

\begin{table}[t]
\centering
\small
\caption{Best ordered triples. The third step is compared against the pair it extends.}
\label{tab:triples}
\begin{tabular}{@{}clcccc@{}}
\toprule
Rank & Chain & Removal $\uparrow$ & Residual TPR@1\%FPR $\downarrow$ & COMET $\uparrow$ & Length ratio \\
\midrule
1 & \ph{chain} & \ph{run} & \ph{run} & \ph{run} & \ph{run} \\
2 & \ph{chain} & \ph{run} & \ph{run} & \ph{run} & \ph{run} \\
3 & \ph{chain} & \ph{run} & \ph{run} & \ph{run} & \ph{run} \\
4 & \ph{chain} & \ph{run} & \ph{run} & \ph{run} & \ph{run} \\
5 & \ph{chain} & \ph{run} & \ph{run} & \ph{run} & \ph{run} \\
\bottomrule
\end{tabular}
\end{table}

\subsection{Frontier}

Figure~\ref{fig:pareto} plots removal against preservation for every evaluated chain. The frontier is the set of chains not dominated on both axes. The recommendation in \S\ref{sec:discussion} is the frontier point with the highest removal among chains at or above $\tau$.

\begin{figure}[t]
\centering
\fbox{\parbox[c][0.28\textheight][c]{0.9\linewidth}{\centering \ph{Pareto plot: removal rate on the vertical axis, the selecting composite $S$ (BERTScore and COMET) on the horizontal axis; single tactics, pairs, and triples in separate series; the semantic floor drawn as a vertical line.}}}
\caption{Removal against preservation for all evaluated chains.}
\label{fig:pareto}
\end{figure}

\subsection{Second-Attacker Ablation}

The top five chains, the best single tactic, and the control are re-run end to end with Claude Sonnet on the held-out split. Table~\ref{tab:ablation} reports removal under both attackers and the change. The rank correlation across the seven rows is $\rho = \ph{value}$, and the largest absolute removal change is \ph{value}. \ph{One sentence on whether the ranking transfers and what the exception is.}

\subsection{Per-Domain Results}

Table~\ref{tab:perdomain} repeats the recommended chain per domain. Domain effects are expected: news writing constrains entities and numbers, question answering constrains facts, and creative writing leaves more room for paraphrase.

\begin{table}[t]
\centering
\small
\caption{Recommended chain per domain at the same floor $\tau$.}
\label{tab:perdomain}
\begin{tabular}{@{}llcc@{}}
\toprule
Domain & Recommended chain & Removal $\uparrow$ & COMET $\uparrow$ \\
\midrule
News writing              & \ph{chain} & \ph{run} & \ph{run} \\
Long-form QA              & \ph{chain} & \ph{run} & \ph{run} \\
Creative writing          & \ph{chain} & \ph{run} & \ph{run} \\
MarkMyWords tasks         & \ph{chain} & \ph{run} & \ph{run} \\
\bottomrule
\end{tabular}
\end{table}

\section{Discussion}
\label{sec:discussion}

\paragraph{What a guarantee can mean.}
The phrase \emph{removed to a good degree of guarantee} can describe three different claims, and the difference matters for how a result should be read. A \emph{statistical} claim says that the same-config test fails to reject the unwatermarked null at the chosen level, on this document, under this key. An \emph{operational} claim says that a detector tuned to a fixed false-positive rate no longer separates attacked from unattacked text. A \emph{scheme-independent} claim says that no detector for the scheme, present or future, could recover the mark. The third claim is not available for any tactic in our catalogue, and the impossibility result of \citet{sand} explains why: a mark that occupies a small share of the high-quality output space can be walked away from by any attacker with a quality oracle. Our results support the first two claims, and we state them in those terms.

\paragraph{The value of composition.}
\ph{One paragraph once the search finishes: does the best pair add removal over the best single, does the best triple add removal over the best pair, and at what preservation cost. State the marginal removal per point of preservation.}

\paragraph{Order effects.}
Because chains are ordered, the same two tactics can produce different outcomes in the two orders. We expect the asymmetry to follow from what each step leaves behind: a token edit applied before a full-stream replacement is overwritten, while the same edit applied after the replacement is the last change the detector sees. Section~\ref{sec:search} scores both orders on the search split and does not extend an order that misses the minimum detectable difference. \ph{Report the measured asymmetry for the top pairs and say whether the direction is consistent across domains.}

\paragraph{Comparison with prior work.}
Reported removal rates in the literature are not directly comparable across papers, because scheme, length, detector threshold, and attacker model all differ. Our numbers agree in direction with \citet{kgwreliability} and \citet{dipper} on paraphrase, with \citet{cwra} on pivot round trips, and with \citet{markmywords} on the relative weakness of simple token edits. \ph{State the quantitative comparisons where the setups overlap, and the reason for any disagreement.}

\paragraph{Practical reading.}
For a practitioner who wants the strongest removal at a fixed fidelity cost, the recommendation is the frontier point in \S\ref{sec:results}. For a scheme designer, the useful quantity is the cost at which the frontier turns: the amount of rewriting needed to push detection below the operating point. \ph{Fill in the turning point and the corresponding per-document edit budget.}

\paragraph{Threat model and scope.}
The study assumes an attacker with black-box access to a generator and to open attack models, and with knowledge of the scheme class. It does not assume the key, a vendor detector, or an ability to query the detector during the attack. These assumptions are the common ones in the removal literature, and results obtained under them do not transfer automatically to an attacker who holds the key or who can iterate against a live detector.

\section{Limitations and Ethics}
\label{sec:limitations}

\paragraph{The oracle is same-config.}
Every removal number in this paper comes from a detector configured with the same scheme parameters and key family used at generation. That is the right choice for a reproducible study, and it is also its main limit. A vendor detector may implement a different test, may aggregate evidence across a longer span, or may combine the statistical signal with retrieval against stored generations, as \citet{dipper} propose. Our results therefore bound the behaviour of the scheme, and they do not certify what a deployed checker will report. We avoid the phrase \emph{watermark-free} for that reason and report removal relative to a named detector.

\paragraph{Semantic metrics are proxies.}
BERTScore, BLEURT, COMET, and MAUVE each capture part of meaning and miss part of it. They are insensitive to some factual corruptions, they reward paraphrase that preserves entities and numbers, and they can be gamed by a tactic that keeps the opening sentence intact and rewrites the rest. The human rating sample is the check on the automatic scores, and it is small. A reader who cares about a specific failure mode, such as a changed number or a dropped attribution, should look at the per-metric table rather than the composite.

\paragraph{Composition compounds loss.}
A chain that satisfies the floor on average can still fail badly on individual documents. We report the distribution of per-document scores, and we flag chains whose lower tail falls below the floor even when their mean does not. The same applies to length: a tactic that shortens text reduces the evidence available to the detector, and shortening is itself a change in content that the composite score may not penalise enough.

\paragraph{Length effects.}
Detector power grows with the number of tokens. A removal rate measured at \ph{length} tokens does not transfer to a paragraph. We report results by length bucket in Appendix~\ref{app:perdomain}, and we do not extrapolate beyond the buckets we measured.

\paragraph{Model and API drift.}
Hosted models change without notice, and a rewrite produced by \ph{hosted model} in \ph{month} may not be reproducible later. We record model identifiers and dates for every call, pin open-weight checkpoints by revision, and treat a hosted-model change as a new experimental condition. The open-weight row exists partly to keep a reproducible anchor.

\paragraph{Language coverage.}
The evaluation is English-centric, with pivot languages chosen for translation quality. \citet{cwra} show that watermark strength varies across languages, so a chain that works in English may behave differently when the source is not English. We treat cross-lingual generalisation as future work.

\paragraph{Ethics.}
The work measures how robust a deployed provenance mark is, which is the kind of measurement that scheme designers and standards bodies need. The experiments run on prompts and texts the authors are authorised to process, and on public benchmark data. We do not evaluate removal on third-party content, we do not build a service, and we report results at a level that supports defence: which tactic class is cheapest to defend against, and where the frontier turns. Findings that a deployed scheme is weaker than reported are shared with the scheme's maintainers before publication.

\section{Conclusion}
\label{sec:conclusion}

Removal of a sampling-time text watermark is a trade-off between how much of the signal survives and how much of the meaning survives. This paper turns the trade-off into a constrained search: a catalogue of twelve tactics and one control, ordered chains of up to three, a semantic floor fixed in advance, and removal measured at both a per-document threshold and a corpus operating point. The result is a frontier rather than a single number, and a recommendation that names its cost.

The method is deliberately boring. Every tactic has one interface, every chain is scored on the same documents, and every headline number is held out from the search that produced it. We hope that structure outlives the specific ranking, which will move as generators and attackers change.

\paragraph{Future work.}
Three directions follow. The first is cross-scheme transfer: whether the ranking of chains holds for green-list, exponential, and distortion-free schemes, or whether it is specific to SynthID-Text. The second is adaptive attack search, where the attacker queries a detector during the search and the detector is allowed to aggregate over longer spans. The third is defence: given the frontier in \S\ref{sec:results}, which watermark design and which detector-side aggregation push the turning point furthest out at the least quality cost.


\bibliographystyle{plainnat}
\bibliography{refs}

\appendix
\section{Prompts and Rating Protocol}
\label{app:prompts}

\subsection{Attack Prompts}

The paraphrase and style tactics use the templates below. Placeholders in angle brackets are filled from the configuration file. The templates are recorded here because a paraphrase attack is only reproducible if its instruction is.

\begin{quote}
\small
\textbf{Instruction paraphrase.} Rewrite the following passage so that the meaning, entities, numbers, and order of ideas are preserved. Change the wording and sentence structure. Return only the rewritten passage.\\
\textbf{Passage.} \texttt{<source>}
\end{quote}

\begin{quote}
\small
\textbf{Discourse paraphrase.} You are rewriting a passage from a longer document. Preserve every factual claim. You may reorder sentences and merge or split them. Return only the rewritten passage.\\
\textbf{Context.} \texttt{<surrounding paragraphs>}\\
\textbf{Passage.} \texttt{<source>}
\end{quote}

\begin{quote}
\small
\textbf{Style normalisation.} Rewrite the passage in a plain, human register: contractions where natural, no em dashes, no stock transitions, no filler. Preserve meaning and length. Return only the rewritten passage.\\
\textbf{Passage.} \texttt{<source>}
\end{quote}

\begin{quote}
\small
\textbf{Back-translation control.} Translate the passage into \texttt{<pivot language>}. Return only the translation.
\end{quote}

\begin{quote}
\small
\textbf{Cross-lingual generation.} \texttt{<prompt translated into pivot language>}
\end{quote}

\subsection{Human Rating Protocol}

Raters see the source and the attacked document side by side, in random order, without knowing which chain produced the document. They answer two questions on a five-point scale: whether the attacked document says the same thing as the source, and whether it reads fluently. A third question asks them to mark any changed fact, number, or attribution. Each document is rated by \ph{number} raters, and agreement is reported as Krippendorff's $\alpha$. Raters are paid \ph{rate}.

\section{Hyperparameters}
\label{app:hyper}

Table~\ref{tab:hyper} records the settings used for the main sweep. The configuration files in \texttt{configs/} are authoritative; this table is a convenience for readers.

\begin{table}[h]
\centering
\small
\caption{Per-tactic settings. Intensity is normalised where the implementation allows it.}
\label{tab:hyper}
\begin{tabular}{@{}llll@{}}
\toprule
Tactic & Intensity parameter & Default & Swept values \\
\midrule
Instruction paraphrase      & Target rewrite ratio      & \ph{value} & \ph{values} \\
Discourse paraphrase        & Lexical diversity         & \ph{value} & \ph{values} \\
Masked-LM infilling         & Fraction of spans masked  & \ph{value} & \ph{values} \\
Style normalisation         & Rule set                  & \ph{value} & \ph{values} \\
Document rewrite            & Chunk size and context    & \ph{value} & \ph{values} \\
Single-pivot back-translation & Pivot language          & \ph{value} & \ph{values} \\
Multi-hop back-translation  & Number of pivots          & \ph{value} & \ph{values} \\
Cross-lingual generation    & Pivot language            & \ph{value} & \ph{values} \\
Synonym substitution        & Fraction of words         & \ph{value} & \ph{values} \\
Random word edits           & Edit rate                 & \ph{value} & \ph{values} \\
Typo and noise injection    & Character error rate      & \ph{value} & \ph{values} \\
Non-origin regeneration     & Decoding settings         & \ph{value} & \ph{values} \\
\bottomrule
\end{tabular}
\end{table}

\paragraph{Search settings.}
Beam width $B = \ph{value}$; semantic tolerance $\epsilon = \ph{value}$; minimum detectable difference for pruning $\ph{value}$; search split fraction $\ph{value}$; repetitions per stochastic tactic $\ph{value}$.

\paragraph{Detector settings.}
Scheme parameters are those used at generation: \ph{list}. Detection threshold $\alpha = 0.05$ per document, and the corpus operating point at $1\%$ false-positive rate.

\section{Additional Results}
\label{app:perdomain}

\subsection{Results by Length}

Detector power depends on length, so we report the recommended chain by length bucket.

\begin{table}[h]
\centering
\small
\caption{Removal and preservation by length bucket for the recommended chain.}
\label{tab:length}
\begin{tabular}{@{}lcccc@{}}
\toprule
Bucket (tokens) & Documents & Removal $\uparrow$ & TPR@1\%FPR $\downarrow$ & COMET $\uparrow$ \\
\midrule
\ph{range} & \ph{count} & \ph{run} & \ph{run} & \ph{run} \\
\ph{range} & \ph{count} & \ph{run} & \ph{run} & \ph{run} \\
\ph{range} & \ph{count} & \ph{run} & \ph{run} & \ph{run} \\
\bottomrule
\end{tabular}
\end{table}

\subsection{Sensitivity to the Floor}

The recommendation is computed at $\tau$, and we repeat it at $\tau - 0.05$ and $\tau + 0.05$ to show how much the choice of floor moves the answer.

\begin{table}[h]
\centering
\small
\caption{Recommended chain at three semantic floors.}
\label{tab:tau}
\begin{tabular}{@{}lccc@{}}
\toprule
Floor & Recommended chain & Removal $\uparrow$ & COMET $\uparrow$ \\
\midrule
$\tau - 0.05$ & \ph{chain} & \ph{run} & \ph{run} \\
$\tau$        & \ph{chain} & \ph{run} & \ph{run} \\
$\tau + 0.05$ & \ph{chain} & \ph{run} & \ph{run} \\
\bottomrule
\end{tabular}
\end{table}

\subsection{Failure Cases}

\ph{Short list of documents where a chain satisfied the floor in expectation but changed a fact, a number, or an attribution. Each entry names the source, the attacked text, and the change.}

\subsection{Second-Attacker Ablation}

The top five chains, the best single tactic, and the control are re-run end to end with Claude Sonnet on the held-out split.

\begin{table}[h]
\centering
\small
\caption{Removal under the search attacker and the second attacker, on the held-out split. $\Delta$ is the second-attacker removal minus the search-attacker removal.}
\label{tab:ablation}
\begin{tabular}{@{}>{\raggedright\arraybackslash}p{5.6cm}ccc@{}}
\toprule
Chain & Removal (Flash) & Removal (Sonnet) & $\Delta$ \\
\midrule
\ph{chain}          & \ph{run} & \ph{run} & \ph{run} \\
\ph{chain}          & \ph{run} & \ph{run} & \ph{run} \\
\ph{chain}          & \ph{run} & \ph{run} & \ph{run} \\
\ph{chain}          & \ph{run} & \ph{run} & \ph{run} \\
\ph{chain}          & \ph{run} & \ph{run} & \ph{run} \\
Best single tactic  & \ph{run} & \ph{run} & \ph{run} \\
Unicode control     & \ph{run} & \ph{run} & \ph{run} \\
\bottomrule
\end{tabular}
\end{table}

\section{Reproducibility Checklist}
\label{app:checklist}

We answer the Machine Learning Reproducibility Checklist v2.0 \citep{pineau} for the protocol at this commit. Items marked \ph{TODO} are pending the runs. Until the result files exist, every "yes" below describes the harness and the preregistered design, and the results section reports the protocol rather than a measurement.

\paragraph{Models and algorithms.}
\begin{itemize}
  \item Clear description of the mathematical setting: yes, \S\ref{sec:problem}.
  \item Clear explanation of assumptions: yes, \S\ref{sec:problem} and \S\ref{sec:limitations}.
  \item Complexity analysis: partial; wall-clock cost by tactic is reported in \S\ref{sec:setup}.
\end{itemize}

\paragraph{Datasets.}
\begin{itemize}
  \item Statistics and splits: yes, Table~\ref{tab:data} and \S\ref{sec:setup}.
  \item Exclusions and preprocessing: yes, \S\ref{sec:setup}.
  \item Downloadable data: \ph{link}, pending the prompt manifest.
  \item Collection process for new data: self-authored prompts from a documented template, plus the MarkMyWords tasks loaded from their release (Appendix~\ref{app:prompts}).
\end{itemize}

\paragraph{Shared code.}
\begin{itemize}
  \item Dependency specification: yes, \texttt{pyproject.toml} with version ranges.
  \item Evaluation code: partial; the statistics, split, strategy parser, and frontier helpers are in \texttt{efface/evaluate/} and \texttt{efface/search/}. Tactic modules, the detector adapter, and the three-pass search land with the first results.
  \item Pretrained models: named in Table~\ref{tab:attackers}; the open-weight generator checkpoints are named, and revisions are recorded in each run manifest.
  \item README with a table of results and the command that produces each: partial; the README carries the reproduction table with commands marked as planned.
\end{itemize}

\paragraph{Experimental results.}
\begin{itemize}
  \item Hyperparameter ranges and selection method: yes, Appendix~\ref{app:hyper}; the operating intensity is chosen by the pass-1 rule in \S\ref{sec:search}.
  \item Number of evaluation runs: yes, \S\ref{sec:setup}; the search log records every evaluated chain and every skipped extension.
  \item Definition of each reported statistic: yes, \S\ref{sec:metrics}.
  \item Central tendency and variation: yes; prompt-clustered bootstrap intervals for rates, document-level Wilson intervals beside them, mean and standard deviation elsewhere.
  \item Runtime or energy cost: \ph{TODO}.
  \item Computing infrastructure: \ph{TODO}.
\end{itemize}


\end{document}
